% ===================================================================
%  TAFEL Nr. 4 — Reifes Alter und Greisenalter.
%  Traduction 2026 d'apres texto/io, controlee sur texto/fr et
%  texto/en. Consigne : voir texto/de/10-tafel-01.tex.
%
%  DEUX SCENES, DEUX COUCHES, ET LA LIGNE NE PASSE PAS OU LE TABLEAU 1
%  L'AVAIT MISE.
%
%  Le tableau 1 avait separe l'abstrait (latin) du maniable
%  (germanique). Cette page-ci separe autre chose : CE QUE LA FAMILLE
%  EST, et CE QU'ELLE FAIT UN JOUR DE CEREMONIE.
%
%  Ce qu'elle est, tout est allemand, et sans une exception :
%     Braut, Bräutigam, Hochzeit, Schwiegersohn, Schwiegertochter,
%     Schwiegervater, Schwiegermutter, Schwager, Schwägerin,
%     Großmutter, Großvater, Enkel, Enkelin, Witwe, Waise, Vormund,
%     Urgroßvater.
%  Un seul prefixe fait toute la belle-famille — « Schwieger- » —, un
%  seul mot tous les beaux-freres et belles-soeurs — « Schwager »,
%  « Schwägerin ». Aucun de ces mots n'a ete emprunte a personne.
%
%  Ce qu'elle fait ce jour-la est francais d'un bout a l'autre :
%     Veranda, Frack, Serviette, Dessert, Champagner, Likör, Ananas,
%     Toast, Bouquet, Bonbons, Oberkellner, Lackschuhe, Kalender,
%     Rheumatismus, Kompresse.
%  Le repas de noces est un rite importe, et il est arrive avec son
%  vocabulaire, comme le compas et le piano du tableau 1.
%
%  ON N'EMPRUNTE PAS LE NOM DE SA GRAND-MERE ; ON EMPRUNTE LE NOM DU
%  DESSERT QU'ON SERT A SA GRAND-MERE. C'est la meme regle que le
%  tableau 2 avait trouvee sur le corps, poussee d'un cran : elle ne
%  separe pas le proche du lointain, elle separe CE QU'ON A TOUJOURS
%  EU de CE QU'ON A APPRIS A FAIRE.
%
%  ET LA SECONDE SCENE DONNE LA MEME LIGNE UNE TROISIEME FOIS, tracee
%  cette fois entre deux personnes de la meme piece. L'allemand a
%  DEUX mots pour chaque maladie, et le choix depend de qui parle :
%
%     ce que le malade dit          ce que le medecin ecrit
%     -------------------------     -----------------------
%     Lungenentzündung              Pneumonie
%     Zahnschmerzen                 Odontalgie
%     Erkältung                     Katarrh
%     Gliederreißen                 Rheumatismus
%     Ohnmacht                      Synkope
%     Schwindel                     Vertigo
%
%  La colonne ecrit les mots de gauche, parce que c'est la
%  grand-mere qui parle a l'alinea 8 et non son medecin — sauf
%  « Rheumatismus » a l'alinea 10, ou c'est justement le mot que le
%  medecin vient de prononcer et que le vieil homme repete. Le
%  fac-simile ido n'a qu'un mot par maladie ; l'allemand en a deux, et
%  choisir lequel est un choix de PERSONNAGE, non de vocabulaire.
%
%  DEUX INVERSIONS, ET C'EST LA MEME, LA PREMIERE DE LA COLONNE.
%  L'allemand pousse tout complement devant le verbe ou devant le
%  participe, de sorte qu'un complement que l'ido met en dernier
%  remonte : « mit seinem Hund die Büsche ab » rangeait le chien (6)
%  avant les buissons (4), et « von seinem Pudel geführt vor der
%  Apotheke » le caniche (21) avant la pharmacie (16). LE REMEDE EST
%  UNIFORME : rejeter le complement apres le groupe qu'il modifie, en
%  apposition — « die Büsche mit seinem Hund ab », « vor der Apotheke
%  …, von seinem Pudel geführt ». C'est la faute recurrente de cette
%  colonne, comme la chaine du genitif l'etait de l'estonienne.
%
%  Enfin un calque, et il est parfait : les souliers vernis du
%  temoin, alinea 14, sont des « Lackschuhe » — souliers-de-laque.
%  L'anglais dit « patent leather », le cuir brevete ; le francais dit
%  le vernis. Trois langues, trois choses vues sur le meme soulier :
%  le brevet, le vernis, la laque.
% ===================================================================

%%K t04-apar-1 apar
\VUsaut{4.0mm}
\VUtitre{15.0pt}{60}{TAFEL Nr. 4}
\VUsaut{1.6mm}
\VUtitre{11.4pt}{0}{Reifes Alter und Greisenalter.}
\VUsaut{3.2mm}

%%K t04-tit-1 sub
\VUcentre{11.4pt}{0}{\textit{Erste Szene.}}
\VUsaut{1.6mm}
\VUtitre{11.4pt}{0}{Das Hochzeitsessen.}
\VUsaut{2.6mm}

%%K t04-tit-2 sub
\VUpk{10.2pt}{40}{Herbst.}
\VUsaut{3.0mm}

%%K t04-c1-01-1 p
1. --- Die jungen Leute sind groß geworden. Einer von ihnen, Johann, heiratet.
Wir sind beim \VUgras{Hochzeitsessen}, das die Familien der Braut und des
Bräutigams um denselben Tisch versammelt.

%%K t04-c1-02-1 p
2. --- Es ist \VUgras{Herbst}, im Monat \VUgras{September}. Der kalte Wind des
\VUgras{Oktober} und des \VUgras{November} hat dem Land sein grünes Laub noch
nicht genommen. Die Abendluft ist mild und duftet; deshalb wird unter der
\VUgras{Veranda} \textsuperscript{(8)}, die auf den Garten hinausgeht, zu Abend
gegessen.

%%K t04-c1-03-1 p
3. --- In der Ferne, auf dem \VUgras{Hügel} \textsuperscript{(3)}, steht das Haus von
Johanns Eltern, ein elegantes \VUgras{Landhaus} \textsuperscript{(1)}, das im Grün
versteckt liegt; schöne Weinberge, die einen ausgezeichneten Wein geben, bedecken
den Hang. Der Winzer pflegt sie.

%%K t04-c1-04-1 p
4. --- Über den Reben zieht eine Kette \VUgras{Rebhühner} \textsuperscript{(2)}
vorbei, aufgescheucht vom herannahenden \VUgras{Förster} \textsuperscript{(5)}. Er
klopft mit dem \VUgras{Gewehr} unter dem Arm und der
\VUgras{umgehängten Jagdtasche} die \VUgras{Büsche} \textsuperscript{(4)} mit
seinem \VUgras{Hund} \textsuperscript{(6)} ab. Er wacht über das Gut und hält die
Wilderer davon ab, das Wild zu vernichten.

%%K t04-c1-05-1 p
5. --- Draußen an der \VUgras{Veranda} klettert ein \VUgras{Spalierweinstock}, an
dem schwere \VUgras{Trauben} \textsuperscript{(7)} hängen.

%%K t04-c1-06-1 p
6. --- Die schönen Herbstblumen, die den \VUgras{Tisch} \textsuperscript{(16)}
schmücken, sind \VUgras{Chrysanthemen} \textsuperscript{(9)}; der \VUgras{Vater}
\textsuperscript{(10)} der Braut hat sich eine ins Knopfloch seines Fracks gesteckt.

%%K t04-c1-07-1 p
7. --- Das Essen geht zu Ende. Der \VUgras{Diener} \textsuperscript{(11)} fängt an,
die Nachspeisen aufzutragen, und wird gleich die Teile des Gedecks abräumen, die
man nicht mehr braucht --- \VUgras{Löffel} \textsuperscript{(14)},
\VUgras{Gabeln} \textsuperscript{(12)}, \VUgras{Messer} \textsuperscript{(13)},
\VUgras{Schüsseln} \textsuperscript{(15)}.

%%K t04-tit-3 sub
\VUpk{10.2pt}{40}{Die Familie.}
\VUsaut{3.0mm}

%%K t04-c1-08-1 p
8. --- Alle sehen fröhlich aus, und das Lächeln, mit dem die \VUgras{Mutter}
\textsuperscript{(17)} des Bräutigams ihre Kinder ansieht, zeigt, wie glücklich sie
ist.

%%K t04-c1-09-1 p
9. --- Diese Heirat verbindet zwei wohlhabende Familien der Gegend. Johann, der
\VUgras{Ehemann} \textsuperscript{(18)}, ist der \VUgras{Sohn} eines großen
Gutsbesitzers und wird der \VUgras{Schwiegersohn} eines erfolgreichen
Fabrikanten.

%%K t04-c1-10-1 p
10. --- Das \VUgras{Paar} ist jung, und doch sind die beiden schon lange
\VUgras{verlobt}. Die \VUgras{junge Frau} \textsuperscript{(19)}, Martha, sieht
bezaubernd aus in ihrem \VUgras{weißen Kleid} mit den Orangenblüten.

%%K t04-c1-11-1 p
11. --- Ihr \VUgras{Schwiegervater} \textsuperscript{(20)} freut sich über eine so
liebe \VUgras{Schwiegertochter}, und Johanns \VUgras{Schwiegermutter}
\textsuperscript{(21)} lächelt, weil ihre \VUgras{Tochter} so hübsch aussieht.

%%K t04-c1-12-1 p
12. --- Am Ende des Tisches sitzen die übrigen \VUgras{Gäste}
\textsuperscript{(22)}: \VUgras{Verwandte}, \VUgras{Freunde}, \VUgras{Vettern}. Ein
\VUgras{Oberkellner} \textsuperscript{(23)} bedient sie flink, die Serviette über dem
Arm.

%%K t04-tit-4 sub
\VUpk{10.2pt}{40}{Die Freunde.}
\VUsaut{3.0mm}

%%K t04-c1-13-1 p
13. --- Rechts am Tisch sitzt eine \VUgras{junge Dame} \textsuperscript{(24)}, eine
\VUgras{Freundin} der Braut, dann der \VUgras{Bruder} der Braut, ihr
\VUgras{Trauzeuge} \textsuperscript{(25)}. Er plaudert mit seiner
\VUgras{Brautjungfer} \textsuperscript{(26)}, der Schwester des Bräutigams, die durch
diese Heirat Marthas \VUgras{Schwägerin} wird. Sie ist eine reizende junge Frau.
Sie trägt schönen \VUgras{Schmuck}, eine \VUgras{Perlenkette} und ein goldenes
\VUgras{Armband}.

%%K t04-c1-14-1 p
14. --- Neben ihnen steht ein Herr und hebt sein Glas: er ist ein \VUgras{Freund}
\textsuperscript{(27)} der Familie und einer der \VUgras{Trauzeugen des Bräutigams}.
Er bringt einen \VUgras{Toast} aus, trinkt auf die \VUgras{Gesundheit} des
jungen Paares und wünscht ihm ein langes Glück. Er ist festlich gekleidet: Frack
und \VUgras{Lackschuhe} \textsuperscript{(28)}. Er führt die \VUgras{Großmutter}
\textsuperscript{(29)}, die \VUgras{Witwe} ist.

%%K t04-c1-15-1 p
15. --- Der junge Mann im \VUgras{Frack} \textsuperscript{(31)} mit dem schmalen
schwarzen Schnurrbart ist Johanns \VUgras{Schwager} \textsuperscript{(30)}. Er ist
ein angesehener Ingenieur. Er sitzt neben einer sehr eleganten \VUgras{jungen
Dame} \textsuperscript{(33)}, Marthas \VUgras{älterer Schwester} und Mutter des
lieben kleinen Mädchens, das am Arm seines Vetters kommt, um eine Blume zu bringen
und \VUgras{guten Abend} zu \VUgras{wünschen}. Diese Dame hat sehr schöne
\VUgras{Ohrringe} und eine elegante \VUgras{Frisur}.

%%K t04-c1-16-1 p
16. --- Der kleine Vetter, der in seinem \VUgras{Kittel} \textsuperscript{(36)} und
seiner weiten \VUgras{Kniehose} \textsuperscript{(37)} so komisch aussieht, tut einem
leid. Er ist eine \VUgras{Waise} \textsuperscript{(35)}. Er hat Vater und Mutter sehr
früh verloren. Sein Onkel ist sein \VUgras{Vormund} \textsuperscript{(38)} und ist
ledig geblieben, um das arme Kind großzuziehen. Er ist ein guter Mann. Er hat eine
kleine Glatze und ist sehr kurzsichtig; er trägt einen \VUgras{Zwicker}.

%%K t04-tit-5 sub
\VUsaut{2.4mm}
\VUpk{10.2pt}{40}{Die Nachspeise.}
\VUsaut{3.0mm}

%%K t04-c1-17-1 p
17. --- Hinter den Gästen ist auf einem Tisch mit einer \VUgras{Decke} die
\VUgras{Nachspeise} \textsuperscript{(39)} aufgebaut. In einer großen Schale liegt
Obst: \VUgras{Pfirsiche} \textsuperscript{(40)}, \VUgras{Birnen}
\textsuperscript{(41)}, \VUgras{Äpfel} \textsuperscript{(42)}, \VUgras{Aprikosen}
\textsuperscript{(43)} und \VUgras{Trauben} \textsuperscript{(44)}. Daneben
\VUgras{Ananas} \textsuperscript{(45)}, eine schöne \VUgras{Torte}
\textsuperscript{(46)}, \VUgras{Likörflaschen} \textsuperscript{(48)} und
\VUgras{Champagner} \textsuperscript{(49)}, und dazu Stapel sauberer
\VUgras{Teller} \textsuperscript{(47)}.

%%K t04-c1-18-1 p
18. --- Der \VUgras{Weinkellner} \textsuperscript{(50)} entkorkt eine \VUgras{Flasche}
\textsuperscript{(52)} alten Weins mit einem \VUgras{Korkenzieher}
\textsuperscript{(51)}. Der \VUgras{Korken} \textsuperscript{(53)} sitzt anscheinend
sehr fest. Dieser Weinkellner hat eine \VUgras{Serviette} \textsuperscript{(54)}
über dem Arm.

%%K t04-c1-19-1 p
19. --- Nach dem Essen gehen die Herren in das Rauchzimmer hinüber, und die Damen
machen sich für den Ball fertig.

%%K t04-tit-6 sub
\VUsaut{5.0mm}
\VUcentre{11.4pt}{0}{\textit{Zweite Szene.}}
\VUsaut{1.6mm}
\VUpk{11.4pt}{40}{Großvaters Geburtstag.}
\VUsaut{2.6mm}

%%K t04-tit-7 sub
\VUpk{10.2pt}{40}{Der Besuch.}
\VUsaut{3.0mm}

%%K t04-c2-01-1 p
1. --- Zehn Jahre sind vergangen. Johanns Eltern sind alt geworden, und sie lieben
ihre drei Enkelkinder sehr, zwei \VUgras{Enkel} \textsuperscript{(2)} und eine
\VUgras{Enkelin} \textsuperscript{(3)}. Weil heute \VUgras{Großvaters Geburtstag}
ist, sind die Kinder gekommen, um ihm zu gratulieren.

%%K t04-c2-02-1 p
2. --- Der kleine Peter ist noch sehr klein; er ist erst drei. Er sieht die
\VUgras{Katze} \textsuperscript{(1)} auf sich zukommen. Er möchte ihr schönes,
seidiges \VUgras{Fell} und ihren langen \VUgras{Schwanz} streicheln; dass unter
diesen zierlichen \VUgras{Pfoten} böse spitze Krallen sitzen, die ihn leicht
kratzen könnten, kommt ihm gar nicht in den Sinn.

%%K t04-c2-03-1 p
3. --- Seine Schwester Gabriele, die ihn an der Hand hält, ist viel ernster, und
Marcel mit seinem großen \VUgras{Mantel} \textsuperscript{(4)} und dem ledernen
\VUgras{Gürtel} \textsuperscript{(5)} sieht ganz wie ein kleiner Mann aus, trotz der
kurzen Hose, unter der man seine \VUgras{Strümpfe} \textsuperscript{(6)} sieht.

%%K t04-c2-04-1 p
4. --- Ihr Vater hat seinen dicken \VUgras{Wintermantel} \textsuperscript{(7)} mit
dem \VUgras{Pelzkragen} \textsuperscript{(8)} angezogen, und in seiner
\VUgras{Tasche} \textsuperscript{(9)} steckt ein Schal. Seine Frau trägt ihren
Filzhut mit \VUgras{Federn} \textsuperscript{(10)}, ihre \VUgras{Jacke}
\textsuperscript{(11)}, ihre \VUgras{Pelzboa} \textsuperscript{(12)} und ihren
\VUgras{Muff} \textsuperscript{(13)}.

%%K t04-c2-05-1 p
5. --- Es ist sehr kalt, jetzt wo der Winter da ist. Den ganzen Tag ist
\VUgras{Schnee} \textsuperscript{(19)} gefallen, und die Dächer der Häuser sind ganz
weiß. Mit der \VUgras{Nacht} wird die Kälte noch schärfer. Am Himmel leuchten der
\VUgras{Mond} \textsuperscript{(17)} und die \VUgras{Sterne} \textsuperscript{(18)}
und durchdringen die \VUgras{Dunkelheit}.

%%K t04-tit-8 sub
\VUsaut{4.2mm}
\VUpk{10.2pt}{40}{Krankheit.}
\VUsaut{3.0mm}

%%K t04-c2-06-1 p
6. --- Der arme \VUgras{Blinde} \textsuperscript{(20)}, der vor der
\VUgras{Apotheke} \textsuperscript{(16)} über den Platz geht, von seinem
\VUgras{Pudel} \textsuperscript{(21)} geführt, ist sehr zu bedauern. Justine, das
\VUgras{Dienstmädchen} \textsuperscript{(14)}, bringt aus der Küche eine
\VUgras{Schale} \textsuperscript{(15)} sehr heißen, reichlich gesüßten
\VUgras{Kräutertee}.

%%K t04-c2-07-1 p
7. --- Die \VUgras{Großmutter} \textsuperscript{(23)}, die sich ihre \VUgras{Brille}
\textsuperscript{(26)} vom Tisch holen will, um das \VUgras{Rezept}
\textsuperscript{(27)} zu lesen, das der \VUgras{Arzt} \textsuperscript{(25)} eben
geschrieben hat, steht aus ihrem Sessel auf und stützt sich auf ihren
\VUgras{Stock} \textsuperscript{(24)}.

%%K t04-c2-08-1 p
8. --- Sie hat oft kleine \VUgras{Beschwerden}, \VUgras{Schwindelanfälle},
\VUgras{Erkältungen} und \VUgras{Zahnschmerzen}; ihre Beine sind schwach und ihre
Augen nicht mehr sehr gut. Trotzdem neckt sie gern ihren alten \VUgras{Arzt}, der
ihr immer eine \VUgras{Arznei} \textsuperscript{(28)} verschreiben will. Heute freut
sie sich sehr, ihre junge Familie um sich zu haben.

%%K t04-c2-09-1 p
9. --- Das gut verschlossene Zimmer ist warm und behaglich. Im marmornen
\VUgras{Kamin} \textsuperscript{(29)} brennt ein \VUgras{prächtiges Feuer}
\textsuperscript{(30)}, und im weichen \VUgras{Licht} der eben angezündeten
\VUgras{Lampe} \textsuperscript{(31)} sitzt es sich gut.

%%K t04-tit-9 sub
\VUsaut{4.2mm}
\VUpk{10.2pt}{40}{Der Großvater.}
\VUsaut{3.0mm}

%%K t04-c2-10-1 p
10. --- Sogar der \VUgras{Großvater} \textsuperscript{(33)} hat vergessen, dass er
krank war, dass sein \VUgras{Rheumatismus} ihn an den \VUgras{Sessel}
\textsuperscript{(34)} fesselt und dass er erst gestern \VUgras{Fieber und
Ohnmachtsanfälle} hatte. Er denkt nicht mehr daran, dass er im letzten Winter eine
\VUgras{Lungenentzündung} hatte und dass ein harter, schmerzhafter \VUgras{Husten}
ihm sehr zugesetzt hat.

%%K t04-c2-10-2 p
Und doch hat er sich nur mühsam erholt. Jetzt geht es ihm ein wenig besser. Aber
sein Bein, das er auf ein \VUgras{Kissen} \textsuperscript{(38)} auf einem kleinen
\VUgras{Schemel} \textsuperscript{(39)} legt, tut ihm auch noch weh.

%%K t04-c2-11-1 p
11. --- Wenn er sorgfältig in seinen warmen \VUgras{Schlafrock}
\textsuperscript{(35)} mit den großen perlmutternen \VUgras{Knöpfen}
\textsuperscript{(36)} gewickelt ist und die kleine \VUgras{Waise}
\textsuperscript{(40)} in ihrer weißen \VUgras{Haube}, ihrem blauen \VUgras{Kleid}
\textsuperscript{(42)} und ihrem \VUgras{Umhang} \textsuperscript{(41)} neben ihm
sitzt und ihm die Zeitung vorliest, dann klagt er nicht.

%%K t04-c2-12-1 p
12. --- Jedes Jahr, wenn der \VUgras{Kalender} \textsuperscript{(43)} das
\VUgras{Datum} des ersten \VUgras{Dezember} wiederbringt, wartet er ungeduldig auf
seine \VUgras{Enkelkinder}, weil er weiß, dass sie dieses wichtige \VUgras{Datum}
nicht vergessen werden.

%%K t04-c2-13-1 p
13. --- Er erinnert sich gern an die fernen Tage, als er selbst dem glorreichen
kaiserlichen \VUgras{Marschall} gratulieren ging, dessen \VUgras{Bildnis}
\textsuperscript{(44)} an der Wand hängt und der der \VUgras{Urgroßvater} der Kinder
ist.
\VUsaut{6.0mm}
\VUfilet{22mm}
